\documentclass[varwidth=18cm, border=10pt]{standalone}

%% Modern Package Set
\usepackage{booktabs}  % For professional table rules
\usepackage{array}     % For extended column definitions
\usepackage{microtype} % For improved justification

\begin{document}

\begin{tabular}{@{} p{4.2cm} p{2.0cm} p{10cm} @{}}
    \toprule
    \textbf{Journal} & \textbf{Country} & \textbf{Justification} \\
    \midrule

    International Journal of Modern Physics A (IJMPA) &
    Singapore &
    \textbf{Primary Candidate.} Published the foundational work by Watson and Musielak. This is the natural audience for continuity and engagement with the existing literature regarding the Chiral Dirac Equation. \\
    \addlinespace

    European Physical Journal C (EPJ C) &
    Germany &
    \textbf{Springer Option.} A leading venue for Particles and Fields. It aligns with the requested typographical standards and offers rigorous review for mathematical physics and phenomenology. \\
    \addlinespace

    Physics Letters B \newline (PLB) &
    Netherlands &
    \textbf{European Option.} Specializes in rapid communications. Suitable if the Chiral Seesaw Mechanism is framed as a significant, novel letter-length result regarding CP violation. \\
    \addlinespace

    Revista Mexicana de F\'{i}sica &
    Mexico &
    \textbf{Mexican Option.} The premier physics journal in the region. It maintains a solid reputation for theoretical physics, offering high regional visibility in Latin America. \\
    \addlinespace

    Journal of Mathematical Physics (JMP) &
    USA &
    \textbf{Alternative.} Published by the AIP. Ideal if the final manuscript emphasizes the Bargmann-Wigner generalization and Clifford algebra derivations over the neutrino phenomenology. \\

    \bottomrule
\end{tabular}

\end{document}
